\chapter{Introduction}\label{ch:intro}

University students spend much of their study time inside PDF textbooks and lecture
slides. Reading is comfortable, but it gives little evidence of whether anything has
actually been understood. A chapter can feel clear on the page and still be forgotten
by the time of the examination. Good study techniques, such as structured note-taking,
self-testing and spaced revision, are well known and well supported by research
\cite{roediger2006,cepeda2006,dunlosky2013}, yet doing them by hand takes more time
than most students can give to every topic.

\textbf{RecaLL} (also called ``Recall AI''; the installed product is named
\emph{Recall}, version~0.1.5) was developed to close this gap. It is a local-first
desktop application that takes an academic document and turns it into a guided study
path: a table of contents and topic tree, Cornell-style study notes for each topic,
Socratic questions that probe understanding, flashcards scheduled by the Free Spaced
Repetition Scheduler (FSRS), generated quizzes, and a chat assistant whose answers cite
the uploaded documents. This report documents how the system was planned, specified,
designed, built, tested and released.

\section{Goals and Objectives of the Project}

The main \textbf{goal} of RecaLL is to move students from \emph{passive reading} to
\emph{active recall}. The system should take care of the mechanical work (structuring a
document, drafting notes, writing questions and scheduling revision) so that the
student's time goes to the part that builds understanding: explaining ideas in their
own words, being questioned on them, and revising them at the right moment. It should
do this on the student's own computer, with the student's own choice of AI model, and
without locking their study material into one vendor's cloud.

The project was motivated by four recurring problems in ordinary study:
\begin{itemize}
  \item \textbf{Passive learning.} Reading encourages recognition rather than recall, so
  a topic can seem understood without ever being tested \cite{roediger2006}.
  \item \textbf{Time-consuming note-taking.} Preparing structured notes from a textbook
  takes significant effort, particularly for mathematical and technical subjects.
  \item \textbf{Poor retention.} Without systematic, spaced revision, learned concepts
  fade, and students often notice this only close to an examination \cite{cepeda2006}.
  \item \textbf{Lack of personalised assessment.} Ordinary study material cannot identify
  the specific misunderstanding behind a wrong explanation.
\end{itemize}

The specific objectives below are written in SMART form (specific, measurable,
achievable, relevant and time-bound). All of them were to be met by the final
submission in November~2026. Where a figure has not yet been measured, it is marked for
the testing chapter (\Cref{ch:test}).

\begin{enumerate}
  \item \textbf{Document import.} Accept PDF, Markdown and plain-text files of up to
  50~MB, validate them (including the \texttt{\%PDF} file signature), and turn each one
  into a topic tree and a set of overlapping text chunks (target 1{,}500 characters,
  150-character overlap), reporting progress to the interface while it works. Result: the text of a 930-page textbook is
  extracted in 3.1~s, and a topic is processed with AI in a median of 18.4~s (\Cref{ch:test}).
  \item \textbf{Structured notes.} Generate, for any topic, a Cornell-style note with a
  fixed seven-part structure (core intuition, invariants and definitions, mechanisms,
  worked example, comparisons, exam pitfalls, self-test cues) with mathematics rendered
  by KaTeX, streamed to the interface as it is produced.
  \item \textbf{Socratic assessment.} For any topic or extracted concept, generate
  exactly two diagnostic questions (causal mechanism first, then counterfactual or
  applied scenario), grade the student's free-text answer on a 0--100 mastery scale in
  four bands, and propose 0, 1 or 1--2 targeted flashcards depending on the score.
  \item \textbf{Spaced revision.} Schedule every flashcard with FSRS-6 (\texttt{fsrs}
  library 6.3.2) at a desired retention of 0.9 and a maximum interval of 36{,}500
  days, with the four ratings Again, Hard, Good and Easy \cite{ye2022fsrs,pyfsrs}.
  \item \textbf{Grounded answers.} Answer questions over the student's own documents
  using the five best-matching chunks, stream the answer token by token, and keep only
  citations (\texttt{[S1]}, \texttt{[S2]}, \ldots) that point to chunks that were
  actually retrieved. Target: 0 fabricated citations shown to the user.
  \item \textbf{Self-testing.} Generate multiple-choice and true/false quizzes of 1--50
  questions at three difficulty levels, and score every attempt deterministically on
  the server.
  \item \textbf{Provider choice and key safety.} Support at least four AI providers
  (Gemini, OpenAI, Groq and a local Ollama server), switchable at run time, with BYOK keys
  saved from the AI Providers settings encrypted by AES-256-GCM and passwords hashed with PBKDF2-HMAC-SHA256 at
  600{,}000 iterations.
  \item \textbf{Quality and delivery.} Run the backend test suite, linting, the
  production build and end-to-end browser journeys automatically on every push to
  \texttt{main}, and build a Windows installer automatically from a version tag.
  Result: the \texttt{v0.1.5} installer was published as a GitHub pre-release; test figures are reported in
  \Cref{ch:test} (all 223 backend tests pass; 50.5\% statement coverage).
\end{enumerate}

\section{Scope of the Work}

\paragraph{In scope.}
Local account registration and login (with optional Supabase sign-in); import of PDF,
Markdown and text documents with table-of-contents and heading detection; a PDF reader
with highlights, annotations, ``explain this passage'' and export of an annotated PDF;
AI-generated Cornell notes with a notebook-style editor; Socratic drills per topic and per
concept, with misconception detection and flashcard suggestions; flashcard sets generated
from documents, topics or highlighted passages; FSRS-based review sessions; the
Knowledge Hub chat with source citations; quiz generation, attempts and results; a
dashboard and learning-progress analytics; AI provider settings with bring-your-own-key
(BYOK) credentials; an editor for the system prompts; a token-usage inspection page; and
a packaged Windows desktop installer built from a version tag.

\paragraph{Out of scope.}
Real payment processing;
a hosted multi-user web service or mobile apps; real-time collaboration between students;
built-in OCR in the default installation (scanned or handwritten PDFs are read with the
optional Marker extractor); training or fine-tuning a custom AI model (all models are used through
their providers' APIs or a local Ollama server); signed installers for macOS and Linux
(the release workflow builds for Windows); and cloud synchronisation, which is left as
future work.

\subsection{Current Situation and Context}

Software that supports academic study generally falls into three categories. The first
is the \textbf{PDF reader}, which displays a document and allows highlighting but offers
no way to check understanding. The second is the \textbf{document chat tool}, which
answers questions about an uploaded file. These tools are fast and convenient, but the
student remains the one asking, and a good answer from the tool does not show that the
student could have produced it. The third is the \textbf{flashcard application}, which is
effective for long-term retention through spaced repetition but depends on the quality of
the cards, and writing good cards from a long textbook is slow.

Research on learning explains why this matters. Retrieval practice (testing oneself)
produces better long-term retention than re-reading the same material
\cite{roediger2006}, and spacing revision over time is more effective than massing it
\cite{cepeda2006}. A large review of study techniques rated practice testing and
distributed practice as the two most useful, and rated highlighting and re-reading as
having low utility \cite{dunlosky2013}. Yet highlighting and re-reading are exactly what a
PDF reader encourages. The Cornell note-taking method \cite{pauk2011} adds a structured
summary and cue questions to each page of notes, but it is rarely followed because it
doubles the effort of note-taking.

In practice, a student who wants all of these techniques must combine several tools: a
PDF reader for the text, a chat tool for explanations, a notes application, and a
flashcard application into which cards are typed by hand. Material is copied between them,
the link back to the source page is lost, and the student's own AI key or local model can
rarely be used.

\subsection{Competing Products (Available in Market)}

\Cref{tab:competitors} compares RecaLL with five well-known products that address parts of
the same problem \cite{notebooklm,chatpdf,anki,quizlet,remnote}. The descriptions are
based on each vendor's published information at the time of writing. These products change
quickly, so the comparison reflects their documentation, not a full evaluation.

\begin{table}[htbp]
\centering
\caption{Comparison with existing products}\label{tab:competitors}
\small
\begin{tabularx}{\textwidth}{@{}>{\raggedright\arraybackslash}p{2.0cm} X X X@{}}
\toprule
\textbf{Product} & \textbf{Key Features} & \textbf{Limitations (for our goal)} & \textbf{Our Advantage}\\
\midrule
NotebookLM & Answers from uploaded sources with inline citations; can generate
flashcards and practice quizzes & Hosted service with a fixed model provider; no FSRS
revision schedule; no Cornell notes & Scheduled FSRS revision, Cornell notes, Socratic
grading; runs on the desktop with a model of the student's choice\\
\addlinespace
ChatPDF & Chat with an uploaded PDF; citations that jump to the passage & Focus is
answering the student's questions, not questioning the student; no revision path &
Socratic drills that question the student and turn gaps into flashcards\\
\addlinespace
Anki & Mature open-source flashcards; FSRS available as an optional scheduler since
version 23.10 & Cards are written or imported by the student; no link from a card to its
source page & Cards generated from the source and linked back to the document and page\\
\addlinespace
Quizlet & Large study platform; AI-generated study guides and flashcards; adaptive Learn
mode & Hosted service; data and model choice controlled by the vendor & Local-first
storage; BYOK or a fully local Ollama model\\
\addlinespace
RemNote & Notes connected to spaced repetition; flashcards from notes and PDFs; FSRS
support & Card generation from the student's own notes and PDFs, but no grading of a
student's free-text explanation & Socratic grading with misconception diagnosis and
targeted cards; local-model option\\
\bottomrule
\end{tabularx}
\end{table}

Each of these tools is strong in its own area. RecaLL's contribution is to join the
complete path, from structuring a PDF to note-making, Socratic questioning, flashcard
generation and FSRS revision, in one desktop application with a free choice of AI
provider, including a model that runs entirely on the student's computer.

\section{System Overview}

RecaLL is a desktop application for students. Its user interface is a React~19 and
TypeScript single-page application inside a Tauri~v2 shell. The interface talks over HTTP
to a local FastAPI backend that listens only on \texttt{127.0.0.1}. In the installed
build, the backend is a PyInstaller executable started by Tauri as a ``sidecar''; during
development it is started by \texttt{run.bat} on port 8000 next to the Vite development
server on port 5173.

The backend owns all processing and storage. Documents, notes, flashcards, conversations
and progress are kept in two local SQLite databases in Write-Ahead Logging (WAL) mode
(an application database, \texttt{recall\_saas.db}, and a per-user study database,
\texttt{user\_<id>\_recall.db}), and uploaded files are kept in the user's application-data
folder. Every AI request goes through one provider layer, which can call Gemini, OpenAI,
Groq or OpenRouter in the cloud, or an Ollama server on the same machine. Cloud
synchronisation through Supabase and tracing through Langfuse are optional and are off
unless configured.

A typical session follows the study path in \Cref{fig:overview}. The student imports a
lecture PDF, reads it in the built-in reader, generates a Cornell note for a topic, takes
a Socratic drill, saves the suggested flashcards, and later reviews the cards that FSRS
reports as due. At any point the student can ask the Knowledge Hub a question about their
documents or generate a quiz, and the dashboard and progress page show what has been
studied and what needs attention next.

\figplaceholder[0.92\textwidth]{system_overview.pdf}{System overview of RecaLL: the student,
the desktop application, the local backend, local storage and the optional AI providers}{fig:overview}

The AI capability is the core of the product, but it is used through providers rather than
trained by the team. The system decides \emph{what} to ask the model (prompts from a central
registry, which the student can edit), \emph{which context} to send (selected chunks and
topic text, isolated as untrusted data), and \emph{how to check} the reply (JSON repair,
schema validation and citation checking). These parts are specified in \Cref{sec:ai-req}
and designed in \Cref{ch:arch,ch:design}.

\section{Structure of the Document}

The rest of this report is organised as follows. \Cref{ch:pm} gives the project
management plan: team organisation and contributions, the incremental process model,
risks, constraints, tools, the timeline, the budget and the impact of the project.
\Cref{ch:req} specifies the requirements through stakeholders, twelve use cases with
textual descriptions, activity, class and sequence diagrams, and the security and AI
requirements. \Cref{ch:arch} describes the layered local client--server architecture, the
technology stack, the AI-integration architecture, the CI/CD pipeline and the Git
workflow. \Cref{ch:design} details the component and low-level design, the two database
schemas, the user interface, the API and the test plan. \Cref{ch:impl} explains how the
key modules were implemented. \Cref{ch:test} reports unit, integration, end-to-end,
security and AI-specific testing with the results obtained. \Cref{ch:deploy} covers
packaging, release and go-live, and \Cref{ch:sustain} presents the sustainability plan and
ethical considerations. The appendices hold the repository details, the full test-case
list, screenshots, meeting minutes, additional diagrams and the AI model card.

\section{Terms, Acronyms, and Abbreviations Used}

\Cref{tab:terms} defines the technical terms and abbreviations used in this report.

\begin{small}\setstretch{1.15}
\begin{longtable}{@{}>{\raggedright\arraybackslash}p{3.0cm} >{\raggedright\arraybackslash}p{0.68\textwidth}@{}}
\caption{Terms, acronyms and abbreviations}\label{tab:terms}\\
\toprule
\textbf{Term} & \textbf{Definition}\\
\midrule
\endfirsthead
\toprule
\textbf{Term} & \textbf{Definition}\\
\midrule
\endhead
\bottomrule
\endfoot
AES-GCM & Advanced Encryption Standard in Galois/Counter Mode, an authenticated cipher. RecaLL
uses AES-256-GCM with a 12-byte nonce to encrypt the BYOK keys saved in its key vault \cite{nist80038d}\\
API & Application Programming Interface\\
BYOK & Bring Your Own Key: the student supplies a personal API key for an AI provider;
requests made with it use no RecaLL credits\\
Chunk & A piece of document text (about 1{,}500 characters) stored with its page number and
used as context for AI features\\
CI/CD & Continuous Integration / Continuous Delivery\\
Cornell note & A note format with main notes, cue questions and a summary \cite{pauk2011};
RecaLL's version has seven fixed parts\\
CORS & Cross-Origin Resource Sharing, the browser rule that limits which origins may call
the local backend\\
E2E & End-to-end (test), here run with Playwright against the built interface\\
Embedding & A numeric vector that represents a piece of text for similarity search\\
FSRS & Free Spaced Repetition Scheduler, an algorithm that models each card's difficulty,
stability and retrievability to choose the next review date \cite{ye2022fsrs}; RecaLL uses
FSRS-6 through the \texttt{fsrs} library\\
IDOR & Insecure Direct Object Reference: reading another user's record by changing an ID\\
JWT & JSON Web Token, a signed token used for stateless authentication \cite{rfc7519};
RecaLL signs its tokens with HS256\\
KaTeX & A JavaScript library that renders LaTeX mathematics in the browser\\
LLM & Large Language Model\\
Mastery score & The 0--100 grade given to a Socratic-drill answer: mastered (85--100),
developing (60--84), fragile (35--59) or misconception (0--34)\\
OCR & Optical Character Recognition, needed for scanned or handwritten PDFs\\
Ollama & A program that runs open-weight language models on the local machine\\
PBKDF2 & Password-Based Key Derivation Function 2; RecaLL hashes passwords with
PBKDF2-HMAC-SHA256, 600{,}000 iterations and a 16-byte random salt \cite{owaspPassword}\\
RAG & Retrieval-Augmented Generation: retrieving relevant text and giving it to the model
as context for its answer \cite{lewis2020rag}\\
SemVer & Semantic Versioning (for example \texttt{v0.1.5}), used for release tags\\
Sidecar & A helper executable shipped with and started by the Tauri application; in RecaLL
it is the packaged FastAPI backend\\
Socratic drill & A short session in which the system asks ``why'' and ``what if'' questions
and grades the student's explanation\\
SQLite & An embedded, file-based relational database\\
SSE & Server-Sent Events: a one-way HTTP stream used to send tokens and progress updates
to the interface\\
Tauri & A framework for building small desktop applications with a web interface and a
Rust core\\
UAT & User Acceptance Testing\\
VAPT & Vulnerability Assessment and Penetration Testing\\
WAL & Write-Ahead Logging, an SQLite journal mode that lets readers work while a writer
commits\\
Workspace & The personal data space that owns a user's documents, conversations and
flashcards in the application database\\
\end{longtable}
\end{small}
